\begin{thebibliography}{99}
\bibitem{ABL} E. Assing, V. Blomer and J. Li,
\emph{Uniform Titchmarsh divisor problems},
\href{https://arxiv.org/abs/2005.13915v1}{arXiv:2005.13915v1} (2020).
The correspondence in Appendix~\ref{app:classical} uses this version.

\bibitem{Pascadi} A. Pascadi,
\emph{Large sieve inequalities for exceptional Maass forms and the greatest
prime factor of $n^2+1$}, Forum of Mathematics, Pi \textbf{14} (2026), e8.
\href{https://doi.org/10.1017/fmp.2026.10025}{doi:10.1017/fmp.2026.10025};
\href{https://arxiv.org/abs/2404.04239v3}{arXiv:2404.04239v3}.

\bibitem{Drappeau} S. Drappeau,
\emph{Sums of Kloosterman sums in arithmetic progressions, and the error term
in the dispersion method}, Proceedings of the London Mathematical Society
\textbf{114} (2017), no.~4, 684--732.
\href{https://doi.org/10.1112/plms.12022}{doi:10.1112/plms.12022};
\href{https://arxiv.org/abs/1504.05549v4}{arXiv:1504.05549v4}.
For the numbering used here, see the
\href{https://drappeau.perso.math.cnrs.fr/documents/divtm-siegel.pdf}{author-hosted version}.

\bibitem{Young} M. P. Young,
\emph{Explicit calculations with Eisenstein series},
\href{https://arxiv.org/abs/1710.03624v2}{arXiv:1710.03624v2} (2017).
\href{https://doi.org/10.1016/j.jnt.2018.11.007}{doi:10.1016/j.jnt.2018.11.007}.

\bibitem{MaynardI} J. Maynard,
\emph{Primes in arithmetic progressions to large moduli I: Fixed residue classes},
\href{https://arxiv.org/abs/2006.06572v2}{arXiv:2006.06572v2} (2021).

\bibitem{MaynardII} J. Maynard,
\emph{Primes in arithmetic progressions to large moduli II: Well-factorable estimates},
\href{https://arxiv.org/abs/2006.07088v1}{arXiv:2006.07088v1} (2020).

\bibitem{Wu} J. Wu,
\emph{Chen's double sieve, Goldbach's conjecture and the twin prime problem, 2},
Acta Arithmetica \textbf{131} (2008), 367--387.
\href{https://arxiv.org/abs/0709.3764}{arXiv:0709.3764}.
The present argument derives its own full pointwise weight and paired deletion.

\bibitem{Lichtman} J. D. Lichtman,
\emph{Primes in arithmetic progressions to large moduli, and Goldbach
beyond the square-root barrier}, with an appendix with S. Drappeau,
\href{https://arxiv.org/abs/2309.08522v1}{arXiv:2309.08522v1} (2023).
The unrestricted moving-residue interface is not an input to the present proof.

\bibitem{LiLiu} Jiamin Li and Jianya Liu,
\emph{Theorem $(1+1.9)$ on Goldbach Conjecture},
\href{https://arxiv.org/abs/2606.05224v2}{arXiv:2606.05224v2} (2026).
The comparison in Section~\ref{sec:current-scope} is limited to the
explicitly identified statements in that version.

\bibitem{Suzuki} Yuta Suzuki,
\emph{The Rosser--Iwaniec sieve}, unpublished notes, 14 March 2023, 102 pp.
\href{https://www.suzuyu1729.com/notes/NagoyaSeminar2022/RosserIwaniec.pdf}{Author-hosted PDF}.
Proposition~9.4(i), pp.~22--23; Proposition~11.8, pp.~64--65;
and Lemma~14.7, pp.~97--98.
\end{thebibliography}
